\chapter{Introduction}
\label{chap:introduction}

\section{Motivation}

The GPU industry has grown significantly in the last decade, driven by
applications far beyond its traditional purpose in graphics rendering.
GPUs are now essential in machine learning, one of today's fastest-growing
industries, as well as in scientific computing, data analytics, and
autonomous systems. As GPU architectures have evolved to meet these
demands, they have become more complex, incorporating sophisticated
memory hierarchies, intricate synchronization mechanisms, and advanced
execution models. Precise specification of GPU instruction set
architectures (ISAs) has therefore become more important than ever:
hardware designers, compiler developers, and software engineers all
depend on accurate ISA documentation to build correct and efficient
systems.

Traditionally, ISA specifications have been written as natural-language
documents, often hundreds of pages long, describing instruction
encodings, semantics, memory models, and execution behaviors. Relying
solely on natural-language specifications has significant drawbacks. Such
specifications can be ambiguous, leaving room for multiple
interpretations of critical behaviors. They are difficult to validate
against actual hardware implementations, often leading to subtle
discrepancies between documented and observed behavior. As ISAs have
grown in complexity, these challenges have become more pronounced.

Executable ISA specifications offer a compelling complement to
traditional ISA documentation. An executable specification is a formal,
machine-readable description of an ISA that can be run. When used
alongside handwritten ISA documents, executable specifications provide
unambiguous semantics that clarify ambiguities in the natural-language
descriptions and resolve points where multiple interpretations are
possible. Because they are executable, these specifications can be
tested against hardware implementations, ensuring
consistency between the documented behavior and actual processor
execution. They serve multiple purposes: as precise reference
documentation that eliminates ambiguity, as executable models for
validation and testing, as foundations for formal verification efforts,
and as tools for generating comprehensive test suites. The CPU
architecture community has demonstrated the value of this approach, with
projects such as the RISC-V formal specification in Sail and ARM's
ASL-based specifications showing how executable models, when paired with
traditional documentation, can significantly improve both specification
quality and the broader ecosystem tooling that depends on clear ISA
semantics.

Despite these advances in CPU ISA specification, GPU architectures have
largely remained documented through traditional means. While some GPU
manufacturers have made portions of their specifications publicly
available, no GPU ISA has an executable specification from its provider.
This gap is particularly significant given the growing role of GPUs in
domains where correctness and reliability are paramount. Moreover, the
SIMT (Single Instruction, Multiple Thread) execution model that defines
GPU behavior introduces unique challenges for formal specification,
including thread divergence and reconvergence semantics, hierarchical
memory spaces, and fine-grained synchronization primitives that differ
substantially from conventional CPU models.

\section{Objective and Scope}

The objective of this thesis is to develop an executable specification of
PISA, the virtual instruction set Intel has published for its
GPUs~\cite{pisa-spec}. The specification is written in ASL, so that it
can be read as documentation of the instruction set and run as its
reference. It covers a subset of PISA, chosen to model execution at the sub-group
level. The subset can exchange data with the host, perform integer
arithmetic, communicate across lanes and branch. The specification
defines its own reconvergence model for divergent control flow, which
PISA leaves open.

To run kernels on the specification, a parser translates PISA text into
the ASL code that loads a kernel into the specification's state, and a
host library presents the specification as an OpenCL device, so that a
host program can launch kernels on it as on any other device The
specification is checked with smoke tests for each instruction and with
kernel workloads compiled from OpenCL.

\section{Contributions}

The contributions of this thesis are the following.

\begin{itemize}
\item An executable specification, in ASL, of a subset of PISA: its
  architectural state, the semantics of each instruction, and an
  execution model for a sub-group of 32 work-items, including
  its own reconvergence model for divergent control flow. To our knowledge, it is
  the only executable specification for PISA.
\item The tools that make the specification run: a parser from PISA text
  to ASL and a host library that exposes the specification through the
  OpenCL interface.
\item A test suite for the specification, with smoke tests for each
  instruction of the subset and a set of kernel workloads compiled from
  OpenCL, and the evaluation of the specification against it.
\end{itemize}

\section{Thesis Structure}

The rest of this thesis is organized as follows.

\begin{itemize}
\item Chapter~\ref{chap:background} gives the background: what an ISA
  specification is and what an executable one adds, PISA and its
  execution model, the state of formal and executable specifications for
  GPU ISAs, and ASL, the language the specification is written in.
\item Chapter~\ref{chap:design} presents the specification and the
  system around it: the subset of PISA that is specified, the ASL
  specification itself, with its state, instruction semantics, execution
  model and handling of divergence, and the parser and host library that
  make kernels run on it.
\item Chapter~\ref{chap:evaluation} evaluates the specification with
  smoke tests for each instruction, checked against the PISA document,
  and with kernel workloads compiled from OpenCL, checked against
  expected outputs and examined through their execution traces.
\item Chapter~\ref{chap:conclusion} concludes, states the limitations of the work and lists future work.
\item Appendix~\ref{app:kernels} gives the source of the evaluation
  kernels.
\end{itemize}